\section{Verification: closed-form hierarchy (CV-1--CV-6)}\label{sec:verif}

We verify the method against six closed-form benchmarks, ordered as an ablation of geometric
difficulty. The first four problems, CV-1 through CV-4, are convex, collinear, or superposable: their
closed forms are linear combinations of known elastic kernels, and on them the radial chart of
\S\ref{sec:tmap} reduces to the ambient level set. The last two, CV-5 and CV-6, break that
equivalence. A cusped non-convex superformula and a self-similar fractal force the transition map and
the neural signed-distance function (SDF) apart, and we measure the separation. We report both regimes
with equal weight. The cases in which the chart buys nothing are not a weakness to hide; they fix the
baseline, so that the CV-5 and CV-6 separation cannot be read as a tuned advantage.

Two conventions hold throughout. We write $\gn<0$ for penetration, where $\gn$ is the normal gap of
\S\ref{sec:tmap}. We solve each benchmark on a learned geometry, a trained neural SDF, a trained
radial chart, or, for the elastostatic problems, a \code{ChartDecoder} domain map, then compare it
against the analytical reference of Appendix~A. Because a neural object can be no more accurate than
the geometry it was trained on, the targets below are matched, not exceeded; we neither expect nor
claim machine precision.

We present the benchmarks one at a time. Each subsection states the physical problem, names the
detector property it tests, gives the closed-form acceptance target, describes how we solve it, reports
the measured result, and ends with the one-line lesson for chart versus level set.

\subsection{The equivalence regime: CV-1--CV-4 (convex / superposable)}\label{sec:verif_equiv}

The first four benchmarks are the elastic and packing baselines, and they share one feature: each
reduces to the superposition of known closed forms \citep{johnson1985contact,timoshenko1970elasticity}.
None presents a non-convex, non-collinear, or curved-surface correspondence to compute, so none
exercises the transition-map machinery. On a convex indenter the radial inverse-chart gap reduces to
the SDF gap, and the two representations are interchangeable at the detector. We state this plainly
because it bounds the claim of the paper: the chart's advantage is specific to hard geometry, and these
four problems are not hard geometry. Their role is to certify the chart-FEM solver and to pin down the
floor that the transition map must clear before any claim of advantage.

% code: benchmarks/contact/cv_numerical/cv1_hertz_fem.py::run_hertz_fem
\subsubsection{CV-1 (Hertz line contact)} Two elastic cylinders pressed together along a line form a
contact patch whose half-width grows with the load. This is the canonical convex contact, and it tests
whether the chart-FEM solver and a neural-SDF indenter reproduce the elastic kernel where chart and
level set agree. The closed form for the half-width under plane strain is
%
\begin{equation}\label{eq:hertz_a}
a = 2\sqrt{\frac{P'R}{\pi E^{*}}},\qquad
\frac{1}{E^{*}} = \frac{1-\nu_1^2}{E_1}+\frac{1-\nu_2^2}{E_2},
\end{equation}
%
where $a$ is the contact half-width, $P'$ the load per unit thickness, $R$ the relative radius of the
two surfaces, $E^{*}$ the reduced modulus, and $E_i,\nu_i$ the Young's moduli and Poisson ratios of the
two bodies. Equation~\eqref{eq:hertz_a} predicts two checks: a constant ratio $a/\sqrt{P'}$ at fixed
$E^{*}$, and the scaling $a\propto (E^{*})^{-1/2}$. To exercise both, we drive a chart-FEM block against
a neural-SDF disc indenter under penalty contact and sweep $E^{*}$ over a factor of eight. The measured
ratio holds the closed form to $\sim 1.6\%$ across the sweep. The contact pressure $p_0 = 2P'/\pi a$ is
the half-ellipse identity, not an independent check, so we do not count it. The lesson: the neural-SDF
indenter matches the exact-SDF baseline rather than improving on it, which is the expected outcome on a
convex body where chart and level set coincide.

% code: benchmarks/contact/cv_numerical/cv2_cattaneo_fem.py::run_cattaneo_fem
\subsubsection{CV-2 (Cattaneo--Mindlin partial slip)} A tangential load below the friction limit, applied
to a Hertz patch, splits the contact into a central zone that sticks and an outer annulus that slips.
This is the simplest frictional contact, and it tests whether the friction layer reproduces the
stick/slip partition on a geometry where chart and level set still agree. With a tangential load
$Q<\mu P$ (where $\mu$ is the friction coefficient and $P$ the normal load) on the CV-1 contact, the
two-dimensional line-contact partition obeys
%
\begin{equation}\label{eq:cattaneo}
\frac{c}{a} = \sqrt{\,1-\frac{Q}{\mu P}\,},
\end{equation}
%
the planar reduction of the Mindlin compliance law
\citep{mindlin1949compliance,johnson1985contact}, where $c$ is the radius of the stick zone and $a$ the
contact half-width of CV-1. We measure $c$ against \eqref{eq:cattaneo} on a deep half-plane, since block
shear does not reproduce the edge-singular tangential traction. The stick radius matches to $\sim
5$--$7\%$ at low $Q$, and the error rises to a mean of $\sim 11\%$ as $Q\to\mu P$. Near the friction
limit the slip boundary sweeps across the contact faster than the node spacing can resolve, so the
discrete stick radius coarsens. The lesson: CV-2 inherits the CV-1 contact and adds only a friction
layer, so the transition map stays dormant and chart and level set remain equivalent.

% code: benchmarks/contact/cv_numerical/cv3_brazilian_fem.py::run_brazilian_fem
\subsubsection{CV-3 (Brazilian disc)} A disc compressed across a diameter develops a tensile stress
transverse to the load and a compressive stress along it. This is a pure stress boundary-value problem
with no contact detector engaged, and it tests the chart-FEM stress recovery on a learned disc
geometry. For a disc of diameter $D$ and thickness $t$ under pole load $P$, the center stress state is
%
\begin{equation}\label{eq:brazilian}
\sigma_{xx} = +\frac{2P}{\pi D t}\quad\text{(tensile)},\qquad
\sigma_{yy} = -\frac{6P}{\pi D t}\quad\text{(compressive)},
\end{equation}
%
with $\sigma_{xx}$ the stress transverse to the load and $\sigma_{yy}$ the stress along it. This state
follows from superposing two Flamant load-point fields and a uniform biaxial tension that restores a
traction-free rim \citep{hondros1959evaluation,timoshenko1970elasticity}. The chart-FEM solve
reproduces the center stresses of \eqref{eq:brazilian} to $1.62\%$ in $\sigma_{xx}$ and $0.58\%$ in
$\sigma_{yy}$ against the exact closed form. The constant-strain triangulation leaves a $\sim 1.6\%$
center floor that plateaus under mesh refinement, a discretization error rather than a modeling error,
and we exclude the pole singularity from the comparison. The lesson: with no detector engaged, the chart
contributes only the domain map, and chart and level set are again interchangeable.

% code: benchmarks/contact/cv_numerical/cv4_nine_disc_fem.py::run_nine_disc_fem
\subsubsection{CV-4 (nine-disc unit cell)} A $3\times3$ array of discs under isotropic confinement loads
each interior disc through four equal contacts at its poles. This is a packing problem whose symmetry
collapses it to a superposition of Brazilian solutions, and it tests whether the charts add by
superposition across multiple contacts. No diagonal contacts form, because $2\sqrt2\,R>2R$; by $D_4$
symmetry every disc sees four equal inward loads $N$ at its north, south, east, and west poles, and
reflection parity makes friction inactive. The unit cell is then the superposition of two perpendicular
Brazilian solutions, with the equibiaxial center stress
%
\begin{equation}\label{eq:nine_disc}
\sigma_{xx} = \sigma_{yy} = -\frac{2N}{\pi R t},
\end{equation}
%
where $N$ is the per-pole contact load, $R$ the disc radius, and $t$ the thickness. The chart-FEM unit
cell reproduces \eqref{eq:nine_disc} to $0.15\%$ in $\sigma_{xx}$ and $0.01\%$ in $\sigma_{yy}$, with a
measured in-plane anisotropy of $0.16\%$ and a center shear stress below the solver tolerance. Two
caveats accompany the result: the equal-$N$ collapse is exact only for rigid loading platens, and this
is the per-disc unit cell rather than a full $N$-body explicit-contact solve, so the contacts only
place the load poles. The lesson: the charts add by superposition exactly as in CV-3, and chart and
level set stay equivalent.

Across CV-1--CV-4 the learned geometry reproduces the contact-mechanics closed forms to $\sim
0.15$--$1.6\%$ (CV-1, CV-3, CV-4) and to $\sim 5$--$11\%$ for frictional partial slip (CV-2), with the
radial chart and the level set equivalent throughout. This is the floor the transition map must clear.
It clears the floor here without buying anything, by construction. The next two benchmarks are where it
starts to earn its place.

\subsection{CV-5: the cusped non-convex superformula, the discriminating test}\label{sec:verif_cv5}

CV-5 is the first benchmark the transition map is built for. A rotating Gielis-superformula cam drives a
free non-convex follower, and the cam's lobe tips are cusps. The problem tests whether the chart can
hold a sharp normal where the ambient SDF cannot. The cam boundary is the superformula radius
\citep{gielis2003generic}
%
\begin{equation}\label{eq:superformula}
\rho(\theta) = \left(\,\left|\tfrac{1}{a}\cos\tfrac{m\theta}{4}\right|^{n_2}
+\left|\tfrac{1}{b}\sin\tfrac{m\theta}{4}\right|^{n_3}\,\right)^{-1/n_1},
\end{equation}
%
where $\rho(\theta)$ is the boundary radius at polar angle $\theta$, $m$ sets the number of lobes,
$a,b$ scale the two axes, and $n_1,n_2,n_3$ shape the lobes; for $n_1<1$ the lobe extrema sharpen into
cusps. This geometry attacks both intrinsic weaknesses of the level set at once. The cusps trigger the
network's spectral bias, and the concave valleys between lobes carry a medial axis on which even the
exact Euclidean SDF gradient is only $C^0$. The boundary chart sidesteps both: it is a one-dimensional
function $\rho(\theta)$ on the circle of directions, with no ambient field and no medial axis, so it
stays sharp where the ambient field cannot.

% code: solvers/contact/supershape.py::SuperformulaRho2D
% code: solvers/contact/supershape.py::radial_gap
We compare three detectors that all consume the matched (gap, normal) contract of \S\ref{sec:tmap}: the
analytical chart, a trained Fourier-feature radial chart $\rho_\theta:\Sph^1\to\R^{+}$, and a neural
SDF of the same shape. For the radial chart, the inverse-chart gap and its conservative normal are
%
\begin{equation}\label{eq:radial_gap_cv5}
\gn(\xbf) = \norm{Q^{\top}(\xbf-\bm{c})} - \rho(\dhat),\qquad
\nbf = \widehat{\,Q\,(\dhat - \grad_{\!S}\rho/r)\,},\qquad
\dhat = \frac{Q^{\top}(\xbf-\bm{c})}{r},
\end{equation}
%
where $\xbf$ is the query point, $\bm{c}$ the cam center, $Q$ the rotation to body frame, $r =
\norm{Q^{\top}(\xbf-\bm{c})}$ the radial distance in that frame, $\dhat$ the unit direction to the
query, and $\grad_{\!S}\rho=(I-\dhat\dhat^{\top})\grad_{\dhat}\rho$ the surface gradient of $\rho$
projected onto the tangent direction. The tangential projection is load-bearing: it strips the spurious
radial component from the raw gradient, so the penalty force is exactly the gradient of a potential.

% code: benchmarks/contact/supershape_cam_drive.py::run_cam_drive
\paragraph{Measured separation.} On the cusped superformula the neural SDF reaches a gap error of
$8\times10^{-3}\,L$ ($L$ the cam length scale) with a degraded normal: the median out-of-plane
component is $\lvert n_z\rvert\sim 0.5$, tilting the normal off the two-dimensional surface. This is the
spectral-bias ceiling, reached by design. The neural radial chart reaches a gap error of
$3.8\times10^{-3}\,L$ and a median normal-angle error of $0.42^{\circ}$, roughly halving the gap error
and recovering an in-plane normal. The chart also resolves a kinematic property the level set's single
closest-point projection cannot: a boundary scan evaluates \eqref{eq:radial_gap_cv5} at every $\theta$
and enumerates every disjoint penetrating arc, reporting $\ge 2$ contact arcs on a wedged pose where a
single closest-point projection returns one foot. Driving the rigid-body cam dynamics with the neural
chart then matches the analytical-chart trajectory to $0.04\%$, while conserving linear and angular
momentum to $5\times10^{-15}$ and $1\times10^{-13}$, respectively. The lesson: the chart resolves
geometry the level set smooths, and reproduces the identical contact dynamics where both detectors
apply. The caveats of \S\ref{sec:tmap} carry over. The radial gap is conservative-large, biased by
$1/\cos\alpha$ on steep flanks (with $\alpha$ the surface slope angle), and the test is rigid-body
contact kinematics, not the deformable solver.

\begin{figure}
  \centering
  \includegraphics[width=\linewidth]{supershape_chart_vs_sdf_pub}
  \caption{CV-5, the discriminating test. A rotating Gielis-superformula cam
  \eqref{eq:superformula} with cusped lobes ($n_1<1$). The trained Fourier-feature radial chart
  $\rho_\theta:\Sph^1\to\R^{+}$ (left) reproduces the analytical gap and normal on the lobe tips and
  in the concave valleys, reaching a gap error of $3.8\times10^{-3}\,L$ and a median normal-angle
  error of $0.42^{\circ}$. A neural SDF of the same shape (right) degrades on the cusps (gap error
  $8\times10^{-3}\,L$) and tilts the normal off the surface (median $\lvert n_z\rvert\sim0.5$), the
  intrinsic spectral-bias ceiling. A boundary scan over the chart enumerates $\ge 2$ disjoint contact
  arcs where a single closest-point projection returns one foot.}
  \label{fig:cv5_chart_vs_sdf}
\end{figure}

\subsection{CV-6: the Koch fractal, resolution independence and the SDF ceiling}\label{sec:verif_cv6}

CV-6 mates two Koch snowflakes, a boundary generated by a four-contraction iterated-function system
(IFS, scale $1/3$) evaluated to recursion depth $n$. The benchmark tests what happens to each
representation as geometric detail refines without bound. The snowflake is not star-shaped for $n\ge
2$, so the radial gap of CV-5 no longer applies. Here the chart is the recursive generating rule
itself, with the IFS doubling as a bounding hierarchy for a pruned recursive descent. The
representation costs of chart and level set are asymmetric, and that asymmetry is the point.

% code: solvers/contact/koch.py::inside_query
% code: solvers/contact/koch.py::KochChart
\paragraph{Storage and query cost.} The IFS chart stores the four contraction maps once, in $O(1)$
storage independent of depth, and evaluates the boundary to any $n$ on demand. Its inside/outside
test prunes the recursion to a per-query cost of $O(\text{depth})$, measured at $\approx 21$ visited
nodes and bounded for $n$ up to $12$. A level set pays differently. A uniform SDF grid storing the same
pattern costs $O(9^n)$ cells, and an adaptive narrow-band SDF costs $O(4^n)$ boundary cells; both are
exponential in depth and capped at their finest cell, whereas the chart sets its resolution per query
(Fig.~\ref{fig:koch_cost}). A precomputed SDF is in fact cheaper \emph{per query}, an $O(1)$ lookup, so
the chart does not win on per-query speed, but it cannot represent the geometry past its stored
resolution. The lesson: the chart's defensible wins are storage and resolution independence, not
query speed.

\begin{figure}
  \centering
  \includegraphics[width=\linewidth]{koch_cost_scaling_pub}
  \caption{CV-6 cost decomposition. The recursive IFS chart stores the Koch boundary in $O(1)$ (four
  contraction maps) and detects contact at $O(\text{depth})\approx21$ visited nodes, bounded for $n$ up
  to $12$. A uniform SDF grid costs $O(9^n)$ cells and an adaptive narrow-band SDF $O(4^n)$ boundary
  cells, both exponential in depth and capped at their finest resolution. The chart trades a cheaper
  per-query lookup for unbounded resolution on demand.}
  \label{fig:koch_cost}
\end{figure}

% code: benchmarks/contact/koch_neural_ceiling.py::train_and_measure
\paragraph{The measured SDF refinement ceiling.} We do not argue the level set's limit on a fractal; we
measure it. A single fixed-capacity neural SDF (12{,}801 parameters), trained on the \emph{exact} Koch
signed distance at each fractal level $n=1\ldots5$, resolves the representable level ($n=1$) and then
plateaus as the detail refines below its resolution. The zero-level-set deviation rises from
$1.9\times10^{-2}$ at $n=1$ and saturates around $4$--$8\times10^{-2}$, while the median normal-angle
error jumps from $7.8^{\circ}$ at $n=1$ to $\sim45^{\circ}$ and never recovers, saturating for $n\ge 2$
(Fig.~\ref{fig:koch_ceiling}). Growing the parameter count would shift the ceiling but not remove it: a
fixed-capacity coordinate network cannot keep resolving self-similar detail, which is the intrinsic
spectral bias of the representation. The lesson: the chart stores the generating rule and stays exact at
any finite level, verified to machine precision against the brute-force polyline distance, where the
level set hits a hard ceiling.

\begin{figure}
  \centering
  \includegraphics[width=\linewidth]{koch_neural_ceiling_pub}
  \caption{CV-6, the measured level-set ceiling. A fixed-capacity neural SDF (12{,}801 parameters)
  trained on the exact Koch signed distance at fractal levels $n=1\ldots5$. The zero-level-set
  deviation and the median normal-angle error jump off the $n=1$ floor (normal error $7.8^{\circ}$) and
  plateau at a capacity ceiling ($\to\sim45^{\circ}$), saturating for $n\ge2$: the network cannot keep
  resolving self-similar detail without growing its parameters. The recursive chart stores the IFS rule
  in $O(1)$ and remains exact at any finite level.}
  \label{fig:koch_ceiling}
\end{figure}

\subsection{Verification summary}\label{sec:verif_summary}

% code: postprocessing/plot_numerical_cv_summary.py::main
Table~\ref{tab:cv_summary} collects the closed-form hierarchy. CV-1--CV-4 verify the chart-FEM solver
on learned geometry in the superposition regime, where the chart and the level set are equivalent and
the measured errors track the closed forms to $\sim0.15$--$1.6\%$ (with frictional partial slip at
$\sim5$--$11\%$). CV-5 and CV-6 are where the transition map is exercised: the radial chart resolves a
cusped non-convex boundary the SDF smooths and drives identical contact dynamics, and the recursive IFS
chart resolves a fractal at any depth while a fixed-capacity SDF hits a measured refinement ceiling.
Figure~\ref{fig:cv_summary} plots the per-benchmark L1 errors against the closed forms alongside the
CV-5 chart-over-SDF separation. The distinctive value is not beating a level set on a sphere, where
they are equivalent, but representing geometry, cusps and fractals, that a single neural SDF cannot.

\begin{table}
  \centering
  \caption{Closed-form verification hierarchy CV-1--CV-6. Errors are the L1 deviation of the numerical
  solve on learned geometry from the analytical reference of Appendix~A. ``Chart vs.\ level set'' states
  whether the radial chart and the neural SDF differ on the benchmark; ``T-map'' marks whether the
  transition-map detector is exercised. Sign convention: $\gn<0$ is penetration.}
  \label{tab:cv_summary}
  \small
  \begin{tabular}{@{}llllcc@{}}
    \toprule
    Benchmark & Closed-form target & Measured & Chart vs.\ level set & T-map \\
    \midrule
    CV-1 Hertz line contact      & $a=2\sqrt{P'R/\pi E^{*}}$        & $1.6\%$ ($a$--$E^{*}$)        & equivalent      & no  \\
    CV-2 Cattaneo--Mindlin       & $c/a=\sqrt{1-Q/\mu P}$           & $5$--$11\%$ ($c/a$)           & equivalent      & no  \\
    CV-3 Brazilian disc          & $\pm2P/\pi Dt,\,-6P/\pi Dt$      & $1.62\%/0.58\%$              & equivalent      & no  \\
    CV-4 Nine-disc unit cell     & $-2N/\pi Rt$ (equibiaxial)      & $0.15\%/0.01\%$              & equivalent      & no  \\
    \midrule
    CV-5 SDF (superformula)      & analytic radial gap/normal      & gap $8\times10^{-3}L$, normal degraded & ---    & ---  \\
    CV-5 chart (superformula)    & analytic radial gap/normal      & gap $3.8\times10^{-3}L$, $0.42^{\circ}$; dyn.\ $0.04\%$ & \textbf{chart wins} & \textbf{yes} \\
    CV-6 Koch fractal            & exact Koch SDF, rising level    & SDF ceiling $7.8^{\circ}\!\to\!45^{\circ}$; chart $O(1)$ & \textbf{chart wins} & \textbf{yes} \\
    \bottomrule
  \end{tabular}
\end{table}

\begin{figure}
  \centering
  \includegraphics[width=\linewidth]{numerical_cv_summary_pub}
  \caption{The numerical CV suite in one view. Left: each benchmark's L1 error against the closed form,
  with CV-1/CV-3/CV-4 at $\le\!\sim\!1.6\%$ and CV-2 at $\sim11\%$ on the coarse mesh. Right: the CV-5
  headline. The trained neural radial chart ($3.8\times10^{-3}L$, $0.42^{\circ}$) separates from the
  neural SDF ($8\times10^{-3}L$, degraded normal) on the cusped superformula, the first benchmark on
  which the transition map and the level set are not equivalent.}
  \label{fig:cv_summary}
\end{figure}
